% !TEX root = ../notes_3dprinting.tex
% Chapter 7: Sequenced Printing -- WRITTEN
% Pause mid print, drop a part into an open cavity, let the printer close
% over it. Not to be confused with print-by-object sequential printing.
%%%%%%%%%%%%%%%%%%%%%%%%%%%%%%%%%%%%%%%%%%%%%%%%%%%%%%%%%%%%%%%%%%%%%%

\index{pause}\index{inserts}\index{magnets}
\chapter{Sequenced Printing}
\printchapterversion{07_sequenced}
\newthought{Stop, Insert, Continue.}

\newthought{Some parts are best built in two sittings.} Print until a cavity is
open and reachable, stop, drop something into it, and let the printer close over
the top. The result is a captive part that could not be assembled afterwards:
a nut trapped in a pocket, a magnet buried flush, a bearing, a PCB, a cable
routed through a wall. Nothing about the machine is special here. The whole
technique is one pause in the right place.

\marginnote{\newthought{Not the same as sequential printing.} Bambu Studio also
offers printing several objects one after another, complete part by complete
part. That is a different feature with a different purpose. This chapter is
about pausing inside a single object.}

\newthought{The pause fires before its layer is printed.} That is the only
timing rule, and getting it backwards is the usual mistake. Scrub the preview
slider up until you find the first layer that closes over the cavity, the layer
where the hole disappears, and put the pause there. Bambu Studio emits a single
line of G-code at that point:

\begin{verbatim}
M400 U1     ; finish queued moves, then wait for me
\end{verbatim}

\newthought{Design the pocket loose, and deeper than the insert.} Give roughly
$0.1$ to $0.2$\,mm of clearance in plan so the insert drops in without force,
and the same again in depth so it sits below the surrounding surface. If the
insert stands proud by even one layer height, the nozzle will strike it on the
next pass. At best that knocks the insert out of place; at worst it shifts the
whole part off the plate and takes the print with it.

\newthought{The procedure, once the model is right:}

\begin{enumerate}[itemsep=2pt]
  \item Slice, then switch to the Preview tab.
  \item Drag the vertical layer slider up to the first layer that spans the
        cavity. Confirm on screen that the hole is gone at that layer and still
        open one layer below.
  \item Right-click the slider at that layer and choose \texttt{Add pause}.
  \item Send the job. The pause travels with the file, so re-slicing after a
        change means re-adding it.
  \item When the printer stops, the toolhead parks and the bed stays hot. Leave
        the plate exactly where it is. Do not lift it, do not rotate it.
  \item Seat the insert and press it flush. Check with a straight edge that
        nothing stands above the top surface.
  \item Resume from the printer screen or from Bambu Studio.
  \item Watch the first layer over the insert the same way you watch a first
        layer on the plate.
\end{enumerate}

\marginnote{\newthought{Have the insert in your hand} before the pause arrives.
The part cools while the printer waits. With PLA a minute costs nothing. With
ABS a long pause invites the layer below to contract and the bond across the
pause line to fail, which is exactly where the part will later break.}

\newthought{Two failure modes are specific to magnets.} Check the polarity
before it goes in, because it is not coming out again, and remember that a
magnet will happily reach through a thin floor and grab the steel plate. Leave
enough material underneath, or accept that removing the finished part becomes a
fight.

\marginnote{\newthought{Teaser.} If the pause fires before its layer is printed,
which layer do you select to bury a magnet whose pocket is $3$\,mm deep?}
